% GPT-6 Astra Pro; effort not exposed; /root/replacement; 22 November 2026.
\section{HD-F1-R13: every pure-character triangle presentation requires at least four terms}
\label{sec:followup-rm-three-pure}

No full balanced character
\[
             w=N e^{-B}(1-qA_1)(1-qA_2),
                    \qquad N\ne0,\quad q>0,
 \tag{F1R13.1}
\]
on the fourfold of R07 is a linear combination of at most
three nonzero characters in the pure-character functor
orbits specified in R09. The coefficients may be arbitrary
nonzero rational numbers. Those orbits include the images
of nonzero exponential characters under the stated actual
Fourier--Mukai equivalences, isogeny pullbacks and pushforwards,
line tensors, translations, shifts and duals. A term is
allowed to have rank zero after such a transport.

Consequently an actual finite triangle filtration of an
object with character (F1R13.1), with all factors having
one of these characters, requires at least four terms
after combining proportional characters and deleting zero
coefficients. The same applies to the character presentation
given by actual standard cohomology sheaves. In particular
every actual adjacent two-cohomology object whose two
endpoint characters belong to these orbits is excluded,
independently of its actual coupling and of any finite action.
This assertion is at the whole-character level; it does
not infer an object from a formal presentation.

\paragraph{The exact maps of the pure terms.}
For a nonzero exponential $z=r e^{\ell}$ on a fourfold,
R09 gives
\[
 c_z^2=r e^{\ell}\bigwedge^2\pi_\ell,
 \qquad
 \pi_\ell:HH^1(X)\twoheadrightarrow H^1(\mathcal O_X),
 \quad\pi_\ell(v,\eta)=\eta+\iota_v\ell.
 \tag{F1R13.2}
\]
Multiplication by $r e^{\ell}$ is invertible on the full
HKR space. Thus $c_z^2$ has rank six. Under any specified
actual transport, the exact R08--R09 cohomology and point-group
maps rewrite its degree-two action in the form
\[
 c_z^2=J_z\bigwedge^2\pi_z,
 \qquad \pi_z:U\twoheadrightarrow W_z,
 \quad\dim W_z=4,
 \quad J_z:\bigwedge^2W_z\lhook\joinrel\longrightarrow HH_{-2}(X).
 \tag{F1R13.3}
\]
Here $U=HH^1(X)$; the map $\pi_z$ is the original quotient
composed with the actual invertible point-group differential,
and $J_z$ is the invertible homology transport composed with
the injection in (F1R13.2). The nonzero isogeny degree and
shift or dual signs enter $J_z$ as explicit invertible
scalars or maps, as in R08--R09. In particular every such
term has exactly rank six, including the rank-zero images.

The full balanced map has rank eighteen and kernel $K_w$.
The quotient lemma of R09, after the precise $B$-shear,
gives for every map (F1R13.3)
\[
                     c_z^2(K_w)\ne0.
 \tag{F1R13.4}
\]
Indeed, vanishing would make $\bigwedge^2\pi_z$ annihilate
the two symplectic five-planes. Each restriction of $\pi_z$
to a four-dimensional block would then have rank at most
one, so $\pi_z$ could have rank at most two, contrary to
its displayed rank four. R10 in fact proves that the
left-hand image has dimension at least two, but only
nonvanishing is needed here.

\paragraph{Rank additivity excludes one, two and three terms.}
Suppose
\[
                     w=\sum_{i=1}^{s}a_i z_i,
                   \qquad 1\le s\le3,\quad a_i\ne0.
 \tag{F1R13.5}
\]
The characteristic contraction is linear in the full
character, giving the exact common-domain identity
\[
                  c_w^2=\sum_{i=1}^s a_i c_{z_i}^2.
 \tag{F1R13.6}
\]
For $s\le2$, its rank is at most $6s<18$, already a
contradiction. For $s=3$, put
$V_i=\operatorname{im}(c_{z_i}^2)\subset HH_{-2}(X)$.
Then
\[
 18=\operatorname{rank}c_w^2
       \le\dim(V_1+V_2+V_3)\le\sum_i\dim V_i=18.
 \tag{F1R13.7}
\]
All inequalities are equalities. The addition morphism
$V_1\oplus V_2\oplus V_3\to HH_{-2}(X)$ is therefore
injective. For every $u\in K_w$, equation (F1R13.6)
is a sum of three vectors in independent image spaces
equal to zero. Each $a_i c_{z_i}^2(u)$ vanishes separately.
Thus every $c_{z_i}^2$ annihilates $K_w$, contradicting
(F1R13.4). This proves the exclusion.

\paragraph{Actual filtrations and checking scope.}
For each actual distinguished triangle $F_{j-1}\to F_j\to Q_j$,
the Grothendieck relation is $[F_j]=[F_{j-1}]+[Q_j]$.
Repeated use of these relations, or of the canonical
truncation triangles, produces the precise character sum
(F1R13.5), with shifts contributing their actual signs.
GRR realizes the same relation and the actual functor
transport in cohomology. The proof then uses only
(F1R13.2)--(F1R13.7); no direct-sum decomposition,
cohomology-gap vanishing, equivariant simplicity or
semiregularity is imposed.

This extends the at-least-four-slope part of DEG\_L S-37
from its stated RM-divisor exponential direct-sum setting
to arbitrary rational algebraic slopes, the complete
specified pure-character functor orbits, and actual
coupled triangle presentations. It does not extend S-37's
particular invariant-Ext bound to those couplings.
R12 supplies an independent elementary proof for two
nonzero-rank exponential terms using the actual divisor
square map. Four-term presentations, arbitrary
non-pure endpoints and all other remaining existence
questions are left unchanged by this exclusion.
